\documentclass[10pt,twocolumn,letterpaper]{article}
\usepackage[review]{cvpr}
\usepackage[utf8]{inputenc}
\usepackage{amsmath,amssymb,amsfonts}
\usepackage{bm}
\usepackage{booktabs}
\usepackage{graphicx}
\usepackage{microtype}
\usepackage{xcolor}
\usepackage{tabularx}
\usepackage{multirow}
\definecolor{cvprblue}{rgb}{0.21,0.49,0.74}
\usepackage[pagebackref,breaklinks,colorlinks,allcolors=cvprblue]{hyperref}
\backrefparscanfalse
\def\paperID{*****}
\def\confName{ICCV}
\def\confYear{2027}
\title{\textbf{Compute Dilution in Budgeted Online RGB-D Gaussian Splatting}}
\author{Anonymous ICCV Submission\\Paper ID: *****}
\begin{document}
\maketitle

\begin{abstract}
Online RGB-D 3D Gaussian Splatting (3DGS) spends a fixed per-frame optimization budget while its map keeps growing. We study \textbf{compute dilution}: as the map grows, established Gaussians receive a shrinking share of the updates. Interventions on two TUM and two Replica scenes (150 frames, eight seeds each) support this as a compute-allocation effect. Raising the optimization budget in proportion to map size, with the spawn allowance fixed, improves raw-error selection on all four scenes ($+1.56$ to $+3.73\,\mathrm{dB}$); with free map growth, PSNR rises monotonically with budget and tracks the fraction of mature Gaussians updated per frame (per-seed Spearman $+1.00$ in every scene). The usual remedy, throttling new Gaussians, is scene dependent: coverage throttling improves raw-error selection on the TUM scenes ($+3.33$ and $+1.77\,\mathrm{dB}$) but lowers it on Replica ($-0.68$ and $-0.87\,\mathrm{dB}$). A throttled policy with temporal-influence selection reaches $27.64\,\mathrm{dB}$ on \texttt{fr2\_xyz}, $2.61\,\mathrm{dB}$ above the strongest budgeted baseline and $0.21\,\mathrm{dB}$ below unconstrained optimization, with a $3.8\times$ smaller map, but it trails raw-error selection on 10 of 14 sequences, including all eight Replica scenes. Coverage also changes densification thresholds and hard caps remove capacity on Replica, so map size is not isolated as the cause. The $15\,\mathrm{ms}$ limit is modeled optimization time rather than end-to-end latency (about 10 FPS measured), and seeds measure repeatability on fixed trajectories, not generalization across scenes. Auxiliary analyses (counterfactual Gaussian utility, a learned utility model, contextual re-ranking, and a hole-rate policy router) are in the supplement.
\end{abstract}

\section{Introduction}
3D Gaussian Splatting (3DGS)~\cite{kerbl3Dgaussians} underlies online RGB-D mapping systems such as SplaTAM~\cite{keetha2023splatam}, Point-SLAM~\cite{sandstrom2023pointslam}, MonoGS~\cite{matsuki2024gaussian}, Photo-SLAM~\cite{huang2024photoslam} and RTG-SLAM~\cite{peng2024rtgslam}.
Deployed on robots and headsets, these systems have only milliseconds of optimization per streaming frame.
At the same time, online pipelines add Gaussians in response to rendering error, so the map grows while the per-frame budget stays fixed; in our runs some maps exceed $100\mathrm{k}$ primitives within 150 frames.

We hypothesize \textbf{compute dilution}: a growing map spreads a fixed optimization budget over more Gaussians and over newly spawned ones, so established (mature) Gaussians are updated less often and quality stagnates or degrades.
We test this with interventions: a corrected-pose selector$\times$throttle factorial, optimization-budget controls (a fixed $\times4$ budget, a budget proportional to map size, and a $\times1/\times2/\times4$ dose-response) that hold the spawn allowance fixed, hard map caps, and per-frame logs of how the budget is allocated by Gaussian age, on two TUM and two Replica scenes.
Raising the budget in proportion to map size improves PSNR on all four scenes ($+1.56$ to $+3.73$ dB), and with free growth the fraction of mature Gaussians updated per frame tracks quality in the dose-response.
The standard remedy, throttling spawns, is not universal: it adds $+3.33$ and $+1.77$ dB on the TUM scenes but lowers PSNR on Replica.
Coverage throttling also changes densification thresholds, and hard caps remove capacity on Replica, so we claim a compute-allocation mechanism under the tested protocol, not that map size is the unique cause.

\textbf{Contributions.}
(1) A measurable formulation of compute dilution, with age-resolved allocation statistics (budget share spent on young Gaussians, selection rate of mature Gaussians) logged per frame.
(2) Interventional evidence on four scenes (two TUM, two Replica; eight seeds each): a $2\times3$ selector--throttle factorial, fixed and map-proportional budget controls, a budget dose-response, and map-cap controls, with their limits stated.
(3) A simple coverage-throttled policy that reaches $27.64$ dB on \texttt{fr2\_xyz} within $0.21$ dB of unconstrained optimization at a $3.8\times$ smaller map, together with a scene-by-scene account of where it fails (10 of 14 sequences, all Replica scenes), including a Replica factorial showing that throttling itself lowers PSNR there.
(4) Auxiliary negative and descriptive results (counterfactual single-Gaussian utility, a learned utility model, contextual re-ranking, a hole-rate router), reported in the supplement with frozen artifacts.

\section{Related Work}
\textbf{Online 3DGS SLAM and dense mapping.}
GS-SLAM~\cite{yan2024gsslam} and SplaTAM~\cite{keetha2023splatam} introduced differentiable 3D-Gaussian SLAM, with SplaTAM using silhouette-guided optimization; Gaussian-SLAM~\cite{yugay2023gaussianslam} adds a sub-map design and LoopSplat~\cite{zhu2025loopsplat} adds loop closure by registering Gaussian sub-maps; Point-SLAM~\cite{sandstrom2023pointslam} uses dynamic neural point clouds; MonoGS~\cite{matsuki2024gaussian} and Photo-SLAM~\cite{huang2024photoslam} push toward real-time operation with loop closure; RTG-SLAM~\cite{peng2024rtgslam} adds compact opaque/transparent roles with rule-based stable/unstable selection and is our primary external policy baseline (re-implemented under exactly matched budgets, since native-system comparison would confound renderer, tracker, and scheduler effects).
Unlike these systems, our object of study is not the full SLAM stack but the compute-allocation layer inside it; we evaluate mapping under ground-truth poses with tracking explicitly decoupled (Sec.~\ref{sec:limitations}).
\textbf{Scalable and efficient Gaussian mapping.}
Recent systems address online map scalability through representation and update design. VTGaussian-SLAM~\cite{hu2025vtgs} ties Gaussians to RGB-D views and avoids keeping every primitive learnable, targeting large scenes and memory use. SEGS-SLAM~\cite{wen2025segs} improves structure and appearance consistency, while DyGS-SLAM~\cite{hu2025dygs} focuses on dynamic-scene localization and reconstruction. In parallel, SteepGS~\cite{wang2025steepgs} derives compact density-control updates for 3DGS reconstruction. A recent online-mapping preprint studies geometry-aware initialization and densification under short optimization schedules~\cite{luu2026geometry}.
These methods make clear that map representation, density control, and online optimization are coupled. Our study asks how growth rate and update policy should respond to coverage and observable scene statistics under a fixed per-frame budget and Gaussian representation. We do not claim density control or efficient Gaussian mapping as new in general. The per-scene SOTA numbers in Appendix Supp.\ Tab.~S2 are protocol-mismatched context.
\textbf{Density control for 3DGS.}
Taming 3DGS~\cite{mallick2024taming} uses guided, constructive densification under resource limits; Mini-Splatting~\cite{fang2024minisplatting} targets a constrained Gaussian count; Pixel-GS~\cite{zhang2024pixelgs} reweights the densification gradient by covered pixels; and 3DGS-MCMC~\cite{kheradmand20243dgsmcmc} replaces heuristic densification and pruning with relocation. These methods address offline reconstruction. Our setting differs in that a per-frame optimization budget is shared between existing and newly spawned Gaussians. We do not benchmark against them, which is a limitation of the present comparison (Sec.~\ref{sec:limitations}).
\textbf{Evaluation protocols in dense SLAM.}
Published systems report full-sequence train-view rendering with global refinement at full resolution (Replica $33$--$42\,\mathrm{dB}$ for MonoGS/Point-SLAM/RTG-SLAM; TUM fr2\_xyz $16$--$25\,\mathrm{dB}$ for MonoGS/SplaTAM/Photo-SLAM) plus full-trajectory ATE. Our protocol is deliberately narrower --- budgeted online mapping on fixed horizons --- so cross-paper numbers are contextual range, never superiority claims (Supp.\ Tab.~S2).
\textbf{Budgeted and utility-driven selection.}
Error-only, error$\times$influence, gradient-norm, and temporal-filtered heuristics all trade quality for compute.
We benchmark them under identical modeled budgets with paired Holm-corrected tests rather than reporting best-of-sweep numbers. Their ordering differs between the evaluated TUM and Replica scenes (Sec.~\ref{sec:experiments}).
\textbf{Submodular maximization and knapsack allocation.}
Budgeted set selection is related to monotone submodular maximization, for which approximation guarantees exist under explicit assumptions and appropriate algorithms~\cite{nemhauser1978analysis,sviridenko2004knapsack}. Those assumptions are not established for our measured reconstruction objective, and our fractional-knapsack ranker does not inherit those guarantees. The measured interactions vary by protocol, so we treat pointwise ranks as references and test map-growth control empirically.

\section{Setting and Method}
\textbf{Budgeted mapping.}
Let $G_t = \{g_1, \dots, g_N\}$ be the active map, $g_i = (\mu_i, \Sigma_i, \alpha_i, c_i)$, with streaming observation $(I_t, D_t)$ and pose $T_t$.
Per frame, a scheduler chooses a subset $A_t \subseteq G_t$ to optimize,
\begin{equation}
\max_{A_t \subseteq G_t} \Delta Q(A_t) \quad \text{s.t.} \quad \sum_{i \in A_t} C_i \le B_t,
\end{equation}
where $\Delta Q$ is the reconstruction improvement, approximated at runtime by a selector score, and $C_i(K) = C_i^{\mathrm{base}} + K \cdot C_i^{\mathrm{step}}$ is a decoupled cost model for $K$ microsteps.
The budget $B_t$ is \emph{modeled} optimization time.
\textbf{Measuring dilution.}
We log two age-resolved statistics per frame: the \emph{young budget share}, the fraction of modeled update budget spent on Gaussians younger than ten frames, and the \emph{mature selection rate}, the fraction of Gaussians at least ten frames old that are selected.
Dilution predicts a falling mature selection rate as the map grows.
\textbf{Selection.}
Candidates are scored by error $\times$ influence mass, optionally with temporal EMA filtering ($\beta = 0.90$), and packed by fractional-knapsack value density $s_i / C_i$ under $B_t = 15\,\mathrm{ms}$ with safety factor $1.10$ (effective cap $13.6\,\mathrm{ms}$).
Each selected Gaussian receives $K = 5$ microsteps; the cost model is calibrated so modeled and wall-clock optimization time agree within measurement noise on the RTX 4050 backend.
Exact Gaussian-to-pixel attribution costs $523.9\,\mathrm{ms}$/frame, so we use a center-sampled projection, $\tilde{m}_i = \alpha_i \det(\Sigma_{2D,i})^{1/4}\, e(\lfloor \mu_{2D,i} \rceil)$, that runs in $0.85\,\mathrm{ms}$ ($614.5\times$ faster) at Spearman $\rho = 0.9079$ and Pearson $r = 0.9028$ against exact attribution. The scheduler totals $2.54\,\mathrm{ms}$ per frame.
\textbf{Coverage throttling.}
Coverage $\kappa_t$ is the fraction of valid pixels with transmission $< 0.5$.
The per-frame spawn cap ramps from $1.0$ to $0.20$ as $\kappa_t$ goes from $0.80$ to $0.90$.
A backlog limit exists but was inactive in the headline runs, because its count was only computed when warmup was enabled; the factorial evaluates it independently.
Warmup pre-allocation is disabled and microsteps are fixed at $K = 5$.


\section{Experiments}
\label{sec:experiments}
\textbf{Protocol.} TUM RGB-D~\cite{sturm2012benchmark} and Replica~\cite{straub2019replica} sequences at $320 \times 240$, $B_t = 15\,\mathrm{ms}$, seeds $[42..49]$; two-sided paired Wilcoxon tests, Holm correction and bootstrap CIs over per-seed trajectory summaries; production \texttt{gsplat} backend; ground-truth poses with tracking decoupled (a mapping benchmark, not SLAM ATE).
With eight seeds the smallest attainable two-sided Wilcoxon $p$ is $0.0078$, and with fixed trajectories seeds measure repeatability only, so we report effect sizes with intervals and do not treat seed-level standardized effects as evidence of generalization.
Artifacts: \texttt{results/phase14\_\allowbreak projected\_\allowbreak area\_\allowbreak fix\_\allowbreak clean\_\allowbreak 7policy\_\allowbreak 8seeds/}.
\textbf{Holdout protocol.} Because online mapping optimizes on the evaluated trajectory, train-view PSNR overstates generalization. We hold out every 8th frame render-only (never mapped, never selected): the train--holdout gap is $+0.47\,\mathrm{dB}$ mean over 30f$\times$3 seeds for both OURS and \texttt{error\_only} (\texttt{results/phase14\_holdout/}). This is a same-trajectory proxy, not independent novel-trajectory generalization.

\begin{table*}[t]
\centering
\small
\begin{tabular}{lcccccc}
\toprule
\textbf{Policy} & \textbf{Mean PSNR} & \textbf{Final PSNR} & \textbf{SSIM} & \textbf{FPS} & \textbf{Map} & \textbf{$\Delta$ vs OURS} \\
\midrule
\texttt{no\_op} & 19.06 & 14.88 & 0.6920 & 90.0 & 23,523 & +8.58 $\checkmark$ \\
\texttt{error\_only} & 24.04 & 20.12 & 0.8838 & 9.8 & 73,383 & +3.59 $\checkmark$ \\
\texttt{error\_influence} & 25.03 & 21.83 & 0.9083 & 9.5 & 105,438 & +2.61 $\checkmark$ \\
\texttt{error\_influence\_temporal} & 24.94 & 21.61 & 0.9079 & 9.5 & 106,392 & +2.70 $\checkmark$ \\
\texttt{rtg\_slam\_reimpl} & 23.16 & 20.04 & 0.8669 & 8.9 & 65,960 & +4.48 $\checkmark$ \\
\textbf{OURS} & \textbf{27.64} & \textbf{28.82} & \textbf{0.9231} & \textbf{11.1} & \textbf{23,765} & --- \\
\textit{FULL} & \textit{27.85} & \textit{29.34} & \textit{0.9245} & \textit{9.6} & \textit{90,582} & $-0.21$ $\checkmark$ \\
\bottomrule
\end{tabular}
\caption{\textbf{Main benchmark, TUM \texttt{fr2\_xyz}, 150 frames $\times$ 8 seeds.} All six controls differ significantly from OURS after Holm correction ($p_{\mathrm{Holm}} = 0.0469$). FULL is $0.21\,\mathrm{dB}$ higher. Headroom recovery is $97.6\%$ relative to no-op and $92.5\%$ relative to the strongest budgeted baseline, at $3.8\times$ smaller map than FULL.}
\label{tab:main}
\end{table*}

\begin{table*}[t]
\centering
\footnotesize\setlength{\tabcolsep}{4pt}
\begin{tabular}{lcccccc}
\toprule
& \multicolumn{3}{c}{TUM \texttt{fr2\_xyz}} & \multicolumn{3}{c}{TUM \texttt{fr1\_xyz}} \\
\textbf{Selector / map control} & \textbf{PSNR} & \textbf{Depth L1} & \textbf{Map (k)} & \textbf{PSNR} & \textbf{Depth L1} & \textbf{Map (k)} \\
\midrule
Error, no throttle & 24.05 & .0236 & 80.3 & 21.12 & .0505 & 80.6 \\
Error, coverage & 27.38 & .0148 & 24.4 & 22.89 & .0443 & 37.9 \\
Error, coverage + backlog & 27.36 & .0149 & 24.4 & 22.95 & .0456 & 35.9 \\
Temporal$\times$influence, no throttle & 24.96 & .0165 & 105.5 & 22.06 & .0484 & 89.1 \\
Temporal$\times$influence, coverage & 27.64 & .0130 & 23.8 & 23.48 & .0425 & 38.7 \\
Temporal$\times$influence, coverage + backlog & 27.66 & .0131 & 23.4 & 23.55 & .0450 & 32.5 \\
\midrule
Error, optimization budget $\times4$ & 27.34 & .0142 & 93.8 & 24.19 & .0402 & 79.3 \\
Error, hard map cap 23k & 26.69 & .0215 & 23.0 & 22.30 & .1023 & 23.0 \\
Error, hard map cap 39k & --- & --- & --- & 21.15 & .1046 & 39.0 \\
\bottomrule
\end{tabular}
\caption{\textbf{Corrected-pose factorial and dilution controls, 150 frames $\times$ 8 seeds per scene.} All factorial cells use the same modeled $15\,\mathrm{ms}$ optimization budget. The $\times4$ control uses $60\,\mathrm{ms}$ for optimization while keeping the original $15\,\mathrm{ms}$ spawn allowance. Depth L1 is in meters. Hard caps harm depth on \texttt{fr1\_xyz}; the 39k cap matches the coverage map size there but gives no PSNR gain. Backlog is enabled independently of warmup only in the last throttle level. Per-run traces and paired intervals are in \texttt{results/dilution\_factorial/}.}
\label{tab:factorial}
\end{table*}

\textbf{Main result (Tab.~\ref{tab:main}, Fig.~\ref{fig:deltas}, Fig.~\ref{fig:pareto}).}
OURS reaches $27.64 \pm 0.04\,\mathrm{dB}$: $+2.61$ over the strongest budgeted baseline, $+4.48$ over the RTG-SLAM-policy re-implementation, and $+8.58$ over pass-through. These results evaluate the combined policy on this trajectory.

\textbf{Selector $\times$ throttle factorial (Tab.~\ref{tab:factorial}).} With selector fixed to raw error, coverage adds $+3.33$ $[3.19,3.50]$ dB on \texttt{fr2\_xyz} and $+1.77$ $[1.73,1.81]$ dB on \texttt{fr1\_xyz}. With temporal$\times$influence fixed, its gains are $+2.68$ $[2.64,2.72]$ and $+1.42$ $[1.41,1.44]$ dB. Brackets are paired seed-bootstrap 95\% CIs; all four coverage contrasts have two-sided paired Wilcoxon $p=.0078$. Temporal selection itself adds $+0.26$ dB with coverage on \texttt{fr2\_xyz} and $+0.59$ dB on \texttt{fr1\_xyz}; the interaction is negative ($-0.65$ and $-0.35$ dB), so selector and coverage effects are not additive. The independent backlog option adds $-0.02$/$+0.02$ dB on \texttt{fr2\_xyz} and only $+0.07$/$+0.07$ dB on \texttt{fr1\_xyz} for error/temporal selection; on the latter scene its lower map size accompanies worse depth L1.

\textbf{Dilution controls and direct allocation.} Holding the spawn allowance at its $15\,\mathrm{ms}$ value, giving raw-error selection $60\,\mathrm{ms}$ for optimization raises PSNR by $+3.29$ $[3.18,3.41]$ and $+3.07$ $[3.04,3.11]$ dB on \texttt{fr2\_xyz}/\texttt{fr1\_xyz}, despite maps of $94\mathrm{k}$/$79\mathrm{k}$. Capping the error-only map at 23k adds $+2.64$ $[2.53,2.76]$ and $+1.18$ $[1.15,1.22]$ dB, but doubles depth L1 on \texttt{fr1\_xyz}. A 39k cap matching that scene's coverage map adds only $+0.04$ $[-0.06,0.11]$ dB and also harms depth. Thus a static size limit does not reproduce coverage throttling's effect on \texttt{fr1\_xyz}. For raw-error selection on \texttt{fr2\_xyz}, coverage reduces the share of modeled update budget spent on Gaussians younger than ten frames from $22.5\%$ to $10.7\%$, while the per-frame selection rate among mature Gaussians rises from $35.8\%$ to $82.3\%$. The corresponding mature selection rate on \texttt{fr1\_xyz} rises from $28.6\%$ to $56.3\%$. Allocation shifts and the budget intervention support compute dilution under the tested protocol, but coverage also changes densification thresholds and the optimization trajectory. Map size alone is not identified as the unique cause.

\textbf{Exploratory transfer to two more TUM scenes} (150f$\times$3 seeds, same six factorial cells): coverage adds $+0.34$/$+0.10\,\mathrm{dB}$ to error/temporal selection on \texttt{fr1\_desk}, but depth L1 worsens from $0.0481$/$0.0476$ to $0.0590$/$0.0595\,\mathrm{m}$. On \texttt{fr3\_static}, coverage adds $+0.77$/$+0.38\,\mathrm{dB}$ and slightly improves depth L1. Adding backlog to temporal selection on \texttt{fr1\_desk} loses $0.59\,\mathrm{dB}$. These scene-level differences limit any claim that one growth rule is optimal across sensor and motion conditions; the three-seed comparisons are descriptive (full traces in \texttt{results/dilution\_factorial/}).

\begin{figure*}[tbp]
\centering
\includegraphics[width=0.92\textwidth]{figs/fig7_dilution.pdf}
\caption{\textbf{Map growth and reconstruction quality over time} on TUM \texttt{fr2\_xyz} and \texttt{fr1\_xyz} (150 frames, eight seeds). Lines show the seed mean and shaded bands show $\pm1$ standard deviation. \texttt{error\_only} versus \texttt{error + coverage} holds the selector fixed to show the coverage effect; \texttt{temporal + coverage} adds the combined policy. Paired sequence-level contrasts are reported in Tab.~\ref{tab:factorial}.}
\label{fig:dilution}
\end{figure*}

\textbf{Replica factorial and proportional budget (Tab.~\ref{tab:replica}).}
We ran the same selector$\times$throttle cells on two Replica scenes. Unlike TUM, coverage throttling lowers raw-error PSNR ($-0.68$ and $-0.87$ dB on \texttt{office0} and \texttt{office1}) and shrinks the map ($113$k to $65$k and $84$k to $40$k); backlog throttling is worse still ($-2.55$ and $-2.04$ dB relative to coverage alone, with depth L1 up by $0.12$ and $0.05$\,m). Without throttling, the temporal-influence selector trails raw error by $1.95$ and $2.33$ dB. Budget interventions help: a fixed $\times4$ budget adds $+2.81$ and $+1.56$ dB, and a budget proportional to map size (mean scale $3.67\times$ and $3.87\times$, map size unchanged) adds $+2.88$ and $+1.56$ dB. On TUM the proportional budget adds $+3.73$ and $+2.99$ dB on \texttt{fr2\_xyz} and \texttt{fr1\_xyz} over the unthrottled baseline measured in the same batch ($24.16$ and $21.54$ dB). These batches use the current code revision, whose unthrottled TUM baselines differ from Tab.~\ref{tab:factorial} by $0.11$ and $0.42$ dB, so we compare only within a batch.
\begin{table*}[t]
\centering
\footnotesize\setlength{\tabcolsep}{4pt}
\begin{tabular}{lcccccc}
\toprule
& \multicolumn{3}{c}{Replica \texttt{office0}} & \multicolumn{3}{c}{Replica \texttt{office1}} \\
\textbf{Condition} & \textbf{PSNR} & \textbf{Depth L1} & \textbf{Map (k)} & \textbf{PSNR} & \textbf{Depth L1} & \textbf{Map (k)} \\
\midrule
Error, no throttle & 30.63 & .0403 & 112.7 & 33.84 & .0161 & 84.2 \\
Error, coverage & 29.95 & .0556 & 65.2 & 32.97 & .0273 & 39.9 \\
Error, coverage + backlog & 27.40 & .1720 & 45.7 & 30.93 & .0744 & 25.9 \\
Temporal$\times$influence, no throttle & 28.68 & .0489 & 115.8 & 31.51 & .0202 & 85.8 \\
Temporal$\times$influence, coverage & 28.14 & .0651 & 66.3 & 31.73 & .0329 & 42.0 \\
Temporal$\times$influence, coverage + backlog & 26.18 & .2385 & 45.8 & 29.79 & .1093 & 24.1 \\
\midrule
Error, optimization budget $\times4$ & 33.44 & .0326 & 107.5 & 35.41 & .0106 & 82.8 \\
Error, budget $\propto$ map size & 33.51 & .0327 & 107.5 & 35.41 & .0108 & 82.7 \\
\bottomrule
\end{tabular}
\caption{\textbf{Replica factorial and map-proportional budget, 99 (\texttt{office0}) and 119 (\texttt{office1}) frames $\times$ 8 seeds per scene.} Same cells as Tab.~\ref{tab:factorial}; all runs use the current code revision and a $15\,\mathrm{ms}$ modeled budget unless stated. The proportional cell sets the optimization budget to $15\,\mathrm{ms}\times N_t/N_{\mathrm{ref}}$ (mean scale $3.7$--$3.9\times$) with the spawn allowance fixed at $15\,\mathrm{ms}$. Coverage and backlog throttling lower PSNR on Replica, whereas more optimization budget raises it.}
\label{tab:replica}
\end{table*}

\textbf{Dose-response (Tab.~\ref{tab:dose}).}
With free growth, PSNR rises monotonically with the budget multiplier on all four scenes (for example $24.24$, $26.49$, $27.82$ dB on \texttt{fr2\_xyz}), and the fraction of mature Gaussians updated per frame rises with it; their per-seed Spearman correlation over the three free-growth cells is $+1.00$ in every scene. Over all six cells it stays $+0.92$ and $+0.75$ on the TUM scenes, where a $23$k cap is close to the natural map size, but not on Replica ($+0.20$ and $-0.26$), where the caps ($60$k, $40$k) remove capacity and PSNR falls by $6$--$11$ dB regardless of budget. Under a cap, $\times2$ and $\times4$ coincide once coverage reaches $1.00$, because the budget exceeds what the capped map can use. Update coverage therefore tracks quality wherever capacity is not binding. These are within-scene interventions on raw-error selection. They show that the system is compute limited and that the limit tightens as the map grows, but more compute improves any compute-limited optimizer, so they do not by themselves show that growth causes the limitation; the cap cells address that only on TUM.
\begin{table*}[t]
\centering
\footnotesize\setlength{\tabcolsep}{4pt}
\begin{tabular}{lccccccccc}
\toprule
& \multicolumn{3}{c}{Free map growth} & \multicolumn{3}{c}{Hard map cap} & & \\
\textbf{Scene (cap)} & $\times1$ & $\times2$ & $\times4$ & $\times1$ & $\times2$ & $\times4$ & $\rho_{\mathrm{all}}$ & $\rho_{\mathrm{free}}$ \\
\midrule
\texttt{fr2\_xyz} (23k) & 24.24 (0.38) & 26.49 (0.58) & 27.82 (0.98) & 26.90 (0.73) & 27.85 (1.00) & 27.85 (1.00) & +0.92 & +1.00 \\
\texttt{fr1\_xyz} (23k) & 21.53 (0.29) & 22.78 (0.55) & 24.45 (0.89) & 22.52 (0.65) & 23.63 (1.00) & 23.63 (1.00) & +0.75 & +1.00 \\
\texttt{office0} (60k) & 30.61 (0.17) & 32.25 (0.34) & 33.43 (0.72) & 24.37 (0.24) & 25.25 (0.49) & 26.07 (1.00) & +0.20 & +1.00 \\
\texttt{office1} (40k) & 33.85 (0.18) & 34.88 (0.33) & 35.42 (0.68) & 22.79 (0.38) & 22.85 (0.75) & 22.90 (1.00) & -0.26 & +1.00 \\
\bottomrule
\end{tabular}
\caption{\textbf{Dose-response in optimization budget, 8 seeds.} Raw-error selection, no throttling; entries are mean PSNR (dB) with the mature-update coverage in parentheses, for a budget of $15\times\{1,2,4\}\,\mathrm{ms}$ with the spawn allowance fixed at $15\,\mathrm{ms}$, with free growth or a hard Gaussian cap. $\rho$ is the per-seed Spearman correlation between PSNR and mature-update coverage over the cells (all six, or the three free-growth cells), averaged over seeds. Once coverage reaches $1.00$ extra budget is unused and PSNR saturates. On Replica the caps remove needed capacity (PSNR $6$--$11$\,dB lower at equal budget), so the cross-regime correlation does not hold there.}
\label{tab:dose}
\end{table*}


\textbf{Observed failure cases.} Two TUM scenes show where the combined policy trails raw-error selection.
\textit{Exploration-heavy motion} (\texttt{fr1\_room}, $-1.04\,\mathrm{dB}$) coincides with a smaller final map ($92\mathrm{k}$ vs $97\mathrm{k}$); the available run does not isolate whether throttling or selection causes the gap.
\textit{Rapid rotation} (\texttt{fr1\_rpy}, $-1.13\,\mathrm{dB}$ mean, finals $24.5$ vs $24.1$) is consistent with a lag from temporal smoothing, but this attribution has not been tested independently.
\textbf{Horizon effect (Fig.~\ref{fig:dilution}):} at 30 frames OURS ties \texttt{error\_only} ($-0.01\,\mathrm{dB}$ on \texttt{fr2\_xyz}); the gap appears at longer horizons. The new age-group logs in Tab.~\ref{tab:factorial} link this divergence to a smaller update share for mature Gaussians in the growing map.
Around frame 65 in seed 42, both error-only variants encounter difficult views; the unthrottled run continues spawning roughly 500--900 splats per frame and its PSNR remains low after frame 70, whereas the coverage run spawns roughly 90--130 and recovers by frame 75. The unthrottled loss rises from about $0.07$ at frame 65 to $0.13$ at frame 70 and stays high; opacity declines while the map expands. This is sustained divergence following shared hard views, not evidence for a single isolated gradient explosion. The logs cannot rule out floaters without a spatial geometry audit.


\textbf{Scene dependence (Tab.~\ref{tab:scenes}).}
The combined policy wins the two confirmatory TUM comparisons ($n=8$): \texttt{fr2\_xyz} $+3.59$ dB and \texttt{fr1\_xyz} $+2.38$ dB over raw-error selection, both $p=0.0078$. It trails raw-error selection on all eight Replica scenes ($-1.93$ to $-4.21$ dB; \texttt{office0} and \texttt{office1} at $n=8$, the rest exploratory at $n=3$) and on \texttt{fr1\_room} and \texttt{fr1\_rpy}, so across 14 sequences it loses mean PSNR on 10.
On Replica \texttt{office0}, raw-error selection with a weak throttle ($\kappa=0.98$) reaches $30.17$ dB versus $26.85$ dB for the combined policy (supplement). The Replica factorial (Tab.~\ref{tab:replica}) locates the loss: the unthrottled raw-error cell is the best of the six factorial cells on both scenes, adding the temporal-influence selector costs $1.95$ and $2.33$ dB, and coverage throttling costs $0.68$ and $0.87$ dB for raw-error selection. We cannot tell whether this reflects Replica's rendering, trajectories or depth quality. Compute allocation still matters there, since a map-proportional budget adds $+2.88$ and $+1.56$ dB. A simple hole-rate router that picks the better arm on the two datasets is in the supplement; its within-scene depth interventions did not reverse the ranking, so we do not use it as evidence for any mechanism.
\begin{table*}[t]
\centering
\small
\begin{tabular}{lcccccc}
\toprule
\textbf{Scene} & \textbf{Frames} & \textbf{OURS} & \textbf{\texttt{error\_only}} & \textbf{$\Delta$} & \textbf{Map (k)} & \textbf{$n$} \\
\midrule
fr2\_xyz (TUM) & 150 & 27.64 & 24.04 & $+3.59$ $\checkmark$ & 24 / 73 & 8 \\
fr1\_xyz (TUM) & 150 & 23.47 & 21.09 & $+2.38$ $\checkmark$ & 39 / 81 & 8 \\
fr1\_desk (TUM) & 150 & 22.55 & 22.34 & $+0.21$ & 27 / 56 & 3 \\
fr1\_room (TUM) & 100 & 15.53 & 16.54 & $-1.01$ & 92 / 97 & 3 \\
fr1\_rpy (TUM) & 100 & 19.92 & 21.06 & $-1.14$ & 45 / 62 & 3 \\
fr3\_static (TUM) & 100 & 25.96 & 25.88 & $+0.08$ & 13 / 18 & 3 \\
office0 (Rep.) & 150 & 26.85 & 30.17 & $-3.32$ $\checkmark$ & 79 / 116 & 8 \\
office1 (Rep.) & 100 & 29.25 & 33.46 & $-4.21$ $\checkmark$ & 42 / 74 & 8 \\
office2 (Rep.) & 99 & 22.08 & 24.67 & $-2.59$ & 70 / 117 & 3 \\
office3 (Rep.) & 99 & 22.45 & 24.47 & $-2.02$ & 77 / 150 & 3 \\
office4 (Rep.) & 99 & 23.58 & 25.93 & $-2.36$ & 81 / 141 & 3 \\
room0 (Rep.) & 99 & 20.85 & 22.78 & $-1.93$ & 82 / 150 & 3 \\
room1 (Rep.) & 100 & 20.42 & 24.29 & $-3.87$ & 47 / 59 & 3 \\
room2 (Rep.) & 99 & 22.86 & 25.88 & $-3.03$ & 78 / 146 & 3 \\
\bottomrule
\end{tabular}
\caption{\textbf{Per-scene combined-policy breakdown, 14 sequences.} $\checkmark$ marks the reported significant comparisons at $n=8$; other entries are exploratory. The combined policy wins the two confirmatory TUM comparisons but loses on both confirmatory Replica comparisons; remaining scenes have exploratory sample sizes.}
\label{tab:scenes}
\end{table*}


\textbf{Systems.} Scheduler overhead is $2.54\,\mathrm{ms}$; optimization-budgeted execution runs at about $10\,\mathrm{FPS}$, so this is \emph{not} end-to-end $15\,\mathrm{ms}$ real time.
A separate baseline odometer (dense photometric+depth Gauss-Newton, no keyframes or loop closure) reaches ATE $0.9$/$2.0\,\mathrm{cm}$ on \texttt{fr2\_xyz}/\texttt{fr1\_xyz} over 100 frames (Supp.\ Tab.~S3); it is a sanity check, not a SLAM claim.
AlexNet LPIPS on \texttt{fr2\_xyz} (150f$\times$3) is $0.211$ for OURS versus $0.294$ for \texttt{error\_only}; depth L1 is in Tab.~\ref{tab:factorial}, and we report no mesh-based surface accuracy.
Published-system numbers use different frames, resolution and views and give range context only (Supp.\ Tab.~S2); short-horizon budget-curve and coverage-knee ablations are in Supp.\ Figs.~S1--S2.
\textbf{Reproducibility.} Every table maps to a frozen bundle and a command in \texttt{paper\_repro/}; the headline bundle records per-policy configs, source hashes, artifact hashes, pose protocol and clean-worktree status of each run set. Prior pose-bugged numbers ($19.24\,\mathrm{dB}$) are quarantined in \texttt{results/final\_confirmation/DEPRECATED.md} and never compared.
Environment: RTX 4050 Laptop GPU (6\,GB), PyTorch 2.11+cu128, \texttt{gsplat}~\cite{ye2025gsplat} 1.5.3, Python 3.12; peak allocation stays under $3.5\,\mathrm{GB}$ with the desktop compositor active, which is the intended deployment envelope for budgeted online mapping.

\begin{figure}[tbp]
\centering
\includegraphics[width=\columnwidth]{figs/fig2_main_deltas.pdf}
\caption{\textbf{Paired mean-PSNR differences} on fr2\_xyz (150 frames, $n=8$). Points show OURS $-$ baseline; bars are 95\% paired-bootstrap CIs. Positive favors OURS. FULL is $0.21\,\mathrm{dB}$ higher; all six contrasts pass Holm correction ($p=0.0469$).}
\label{fig:deltas}
\end{figure}

\begin{figure}[tbp]
\centering
\includegraphics[width=\columnwidth]{figs/fig4_pareto.pdf}
\caption{\textbf{Quality--map-size trade-off} on fr2\_xyz (150 frames, $n=8$). Only OURS and FULL lie on the frontier; FULL gains $0.21\,\mathrm{dB}$ with a $3.8\times$ larger map. The three growing heuristics exceed $65\mathrm{k}$; no-op stays compact at lower quality.}
\label{fig:pareto}
\end{figure}

\begin{figure*}[t]
\centering
\includegraphics[width=\textwidth]{figs/fig5_qualitative.pdf}
\caption{\textbf{Qualitative comparison across both evaluated substrates.} Top: an intermediate TUM \texttt{fr2\_xyz} view (90), where OURS has lower per-view error than \texttt{error\_only}; bottom: Replica \texttt{office0} view 98, where the ranking reverses. Columns show ground truth, OURS, \texttt{error\_only}, FULL, and signed RGB $L_1$ error gap (red: OURS has lower error). Insets mark the same crop in each row. PSNR labels are per-view values from seed 42, shown only to illustrate the visual comparison; they are not sequence means.}
\label{fig:qual}
\end{figure*}

\section{Limitations}
\label{sec:limitations}
\textbf{Scope of the mechanism.} The budget interventions cover four scenes but one trajectory per scene and raw-error selection only. Coverage throttling helps on TUM and hurts on Replica, changes both spawn allowance and error thresholds, and map caps neither isolate map size nor remain valid on Replica, where they remove needed capacity. More budget improves any compute-limited optimizer, so the budget interventions show a compute limitation that tightens with map size, not that growth itself is harmful. The new tables use the current code revision and are compared only within a batch (their TUM baselines differ from Tab.~\ref{tab:factorial} by $0.11$ and $0.42$ dB). Additional TUM scenes are exploratory ($n=3$) and show smaller or depth-degrading effects.
\textbf{Protocol.} Mapping uses ground-truth poses; horizons are at most 150 frames (about five seconds at 30 FPS); PSNR is train-view with only a same-trajectory holdout proxy. The $15\,\mathrm{ms}$ budget is modeled optimization time, not frame latency, and the pipeline runs at about 8--11 FPS on a laptop GPU. We do not compare against offline density-control methods~\cite{mallick2024taming,fang2024minisplatting,zhang2024pixelgs,kheradmand20243dgsmcmc} or integrate the policy into a full SLAM system with its own tracker.
\textbf{Statistics.} Eight seeds establish repeatability on fixed trajectories; the unit of generalization is the scene, and only two scenes are confirmatory.
\textbf{Open items.} Longer sequences with a real tracker, comparisons against density-control baselines, additional scenes and sensors, and capacity-matched cap controls on Replica would test whether dilution generalizes beyond these four scenes.

\section{Conclusion}
Under a fixed modeled optimization budget, a growing map leaves established Gaussians under-updated. Raising the budget in proportion to map size improves raw-error selection on all four scenes tested ($+1.56$ to $+3.73\,\mathrm{dB}$), PSNR rises monotonically with budget, and the fraction of mature Gaussians updated per frame tracks quality wherever map capacity is not binding. The remedy of throttling spawns is scene dependent: it recovers several dB on the TUM scenes ($27.64\,\mathrm{dB}$ on \texttt{fr2\_xyz}, $0.21\,\mathrm{dB}$ below unconstrained optimization at a $3.8\times$ smaller map) but lowers PSNR on Replica, where the combined policy trails raw-error selection. We therefore present dilution as a compute-allocation mechanism with scene-dependent remedies, and leave open whether it persists on longer sequences and inside full SLAM systems.

{\small
\bibliographystyle{ieeenat_fullname}
\bibliography{refs}
}

\end{document}
